% Generated draft: requires mathematical and source review.
\documentclass[11pt]{article}
\usepackage[margin=25mm]{geometry}
\usepackage{fontspec}
\setmainfont{Latin Modern Roman}
\usepackage{amsmath,amssymb,mathtools,unicode-math}
\setmathfont{Latin Modern Math}
\usepackage{enumitem,needspace,xurl,hyperref}
\setlength{\parindent}{0pt}
\setlength{\parskip}{5pt}
\setlength{\emergencystretch}{3em}
\newcommand{\N}{\mathbb{N}}
\newcommand{\Z}{\mathbb{Z}}
\newcommand{\Q}{\mathbb{Q}}
\newcommand{\R}{\mathbb{R}}
\newcommand{\C}{\mathbb{C}}
\newcommand{\F}{\mathbb{F}}
\newcommand{\A}{\mathbb{A}}
\newcommand{\PP}{\mathbb{P}}
\newcommand{\E}{\mathbb{E}}
\newcommand{\eps}{\varepsilon}
\newcommand{\dd}{\,\mathrm{d}}
\begin{document}
\section*{UnsolvedMath ID: 20000244}
\textbf{Source problem code:} catalogue entry
\subsection{Definitions and mathematical statement}
Aravind Asok's question asks for which dimensions \(d\) there is a smooth affine \(d\)-dimensional variety over an algebraically closed field carrying a stably free but nonfree projective module of rank \(d-2\), and whether there is a pattern beyond Mohan Kumar's known examples. If \(X=\operatorname{Spec}(A)\), a rank-\(n\) projective module \(P\) is stably free when \(P\oplus A^r\cong A^{n+r}\) for some \(r\geq0\), and nonfree when \(P\not\cong A^n\). The original does not specify the field characteristic or whether existence is required over each algebraically closed field. The extracted text is OCR-damaged and has no clean statement.

The supplied report attributes to Mohan Kumar a published construction of nonfree rank-\(p\) stably free modules on smooth affine \((p+2)\)-folds for prime \(p\), with one-line stabilization. It does not provide the construction's precise ground-field restrictions.

Its candidate prime-divisor propagation theorem is: for every algebraically closed field \(F\) of characteristic zero, every odd integer \(n\geq3\), and every prime divisor \(\ell\mid n\), there exists a smooth affine integral \(F\)-variety \(X\) of dimension \(n+2\) carrying at least \(\ell-1\) pairwise nonisomorphic projective modules \(P_1,\ldots,P_{\ell-1}\) of rank \(n\), such that
\[
P_j\oplus\mathcal O_X\cong\mathcal O_X^{n+1},\qquad
P_j\not\cong\mathcal O_X^n\quad(1\leq j\leq\ell-1),
\]
\[
P_i\not\cong P_j\quad(i\ne j).
\]
Here modules are equivalently their associated vector bundles, and the single \(X\) may depend on \(F,n,\ell\); the source does not assert simultaneous realization for every prime divisor on a fixed \(X\). In dimension notation this would provide examples for every odd \(d\geq5\) with rank \(d-2\). The report says its ingredients are a dimension-flexible use of Mohan Kumar's recursive hypersurfaces, Wendt's odd-rank motivic detector, and a common spreading-out argument, but supplies none of those details.

The report distinguishes the low-dimensional boundary: for \(d\leq3\) the requested nonfree stably free rank is impossible or nonpositive, \(d=4\) is the published \(p=2\) case, and even \(d\geq6\) remain unresolved by this method. This last phrase is a limitation of the proposed method, not an independent claim about the complete literature or every positive characteristic.
\subsection{Short English statement}
Determine which dimensions admit nonfree stably free modules of rank two below the variety dimension. The report cites prime-rank examples and proposes characteristic-zero examples in every odd dimension at least five, with several modules on one variety.
\subsection{Sources}
\begin{itemize}
\item AIM workshop ``Projective modules and \(A^1\)-homotopy theory,'' Question 7, attributed to Aravind Asok: \url{https://aimath.org/pastworkshops/projectiveA1problems.pdf}. No PDF was accessed here.
\item Local canonical registry \texttt{aim-algebraic-geometry-notes.json}, \texttt{notes[243]}, tag \texttt{question}; permanent UnsolvedMath ID \(20000244\), source code \texttt{catalogue entry}. The source names Mohan Kumar's examples and asks whether all \(d\) occur or a pattern exists. The verification field notes visible PDF-extraction damage. The clean statement is absent; original remarks and literature field are empty.
\item Supplied research attempt 1 refers to Mohan Kumar's recursive hypersurfaces and Wendt's odd-rank motivic detector and proposes the prime-divisor propagation theorem. No full authoritative bibliographic references are supplied. The report is machine generated, not peer reviewed, and requires independent verification before citation.
\end{itemize}
\end{document}
